% Casey/W33 phase-defect insertion
% Standalone theorem surface; safe to \input after the graph-RH section.

\part*{Casey--$W(3,3)$ Phase Defect and the Classical Transfer Frontier}
\addcontentsline{toc}{part}{Casey--$W(3,3)$ Phase Defect and the Classical Transfer Frontier}

\section{Completed-$\xi$ regularization}

\begin{definition}[Completed Riemann function]
For $s\in\mathbb C$, let
\[
  \xi(s)=\frac12 s(s-1)\pi^{-s/2}\Gamma(s/2)\zeta(s).
\]
The apparent singularities at $s=0,1$ are removable and
\[
  \xi(0)=\xi(1)=\frac12.
\]
\end{definition}

\begin{remark}
This is the rigorous version of the proposed ``topological anchor.''  No
arbitrary compactification is needed: the completed function itself replaces
the zeta pole by finite analytic endpoint data.
\end{remark}

\section{The exact phase current}

Away from the zero set of $\xi$, define
\[
  \tau(s)=\partial_\sigma\arg\xi(s)
  =\Im\frac{\xi'(s)}{\xi(s)},
  \qquad s=\sigma+it.
\]
Its complete decomposition is
\[
  \frac{\xi'(s)}{\xi(s)}
  =\frac1s+\frac1{s-1}-\frac12\log\pi
   +\frac12\psi(s/2)+\frac{\zeta'(s)}{\zeta(s)}.
\]
Thus the endpoint factors, gamma factor, trivial zeros, and prime-sensitive
zeta term must all be retained.

\begin{proposition}[Argument-principle phase law]
For a positively oriented contour $C$ on which $\xi$ is nonzero,
\[
  \frac{1}{2\pi i}\oint_C\frac{\xi'(s)}{\xi(s)}\,ds
\]
is the number of zeros of $\xi$ inside $C$, counted with multiplicity.
\end{proposition}

\section{Reflected zero orbits}

Let
\[
  \rho=\frac12+\delta+i\gamma,
  \qquad
  \rho^\star=1-\overline\rho=\frac12-\delta+i\gamma.
\]
The functional equation and real symmetry send every zero to this reflected
same-height partner.

The phase derivative of the product
\[
  (s-\rho)(s-\rho^\star)
\]
is
\[
 -\frac{t-\gamma}{(\sigma-\frac12-\delta)^2+(t-\gamma)^2}
 -\frac{t-\gamma}{(\sigma-\frac12+\delta)^2+(t-\gamma)^2}.
\]
The two terms add, so the product current does not vanish when $\delta=0$.
This is the sign obstruction in the original torque derivation.

\section{Reflection cocycle and positive defect energy}

\begin{definition}[Reflection cocycle]
Define
\[
  R_\rho(s)=\frac{s-\rho}{s-\rho^\star}
\]
and its horizontal phase current
\[
  A_\rho(s)=\Im\frac{d}{ds}\log R_\rho(s).
\]
Explicitly,
\[
  A_\rho(\sigma,t)
  =-\frac{t-\gamma}{(\sigma-\frac12-\delta)^2+(t-\gamma)^2}
   +\frac{t-\gamma}{(\sigma-\frac12+\delta)^2+(t-\gamma)^2}.
\]
\end{definition}

\begin{theorem}[Orbit defect energy]
Put $a=\sigma-\tfrac12$ and suppose $a>0$ and $|\delta|<a$.  Then
\[
  \mathcal E_\rho(\sigma)
  :=\int_{-\infty}^{\infty}|A_\rho(\sigma,t)|^2\,dt
  =\frac{\pi\delta^2}{a(a^2-\delta^2)}.
\]
In particular, at $\sigma=1$,
\[
  \mathcal E_\rho(1)
  =\frac{8\pi\delta^2}{1-4\delta^2}.
\]
Hence $\mathcal E_\rho\ge0$, and
\[
  \mathcal E_\rho=0
  \quad\Longleftrightarrow\quad
  \Re(\rho)=\frac12.
\]
\end{theorem}

\begin{proof}
Writing $y=t-\gamma$ gives
\[
  A_\rho=-\frac{4a\delta y}
  {(y^2+(a-\delta)^2)(y^2+(a+\delta)^2)}.
\]
The rational integral evaluates to the stated expression.  Positivity and the
zero criterion are immediate.
\end{proof}

\begin{remark}[Status]
For a finite collection of reflected zero orbits, the sum of the orbit
energies vanishes exactly when every orbit lies on the critical line.  For the
classical $\xi$ zero set, this is an equivalent defect detector, not yet a
proof: one must derive a convergent weighted identity forcing the total energy
to vanish.
\end{remark}

\section{Exact $W(3,3)$ phase operator}

Let $A_\perp$ denote adjacency restricted to the $39$-dimensional nonconstant
sector of the $W(3,3)$ collinearity graph.  Its spectrum is
\[
  2^{24},\qquad (-4)^{15}.
\]
Since $A_\perp/(2\sqrt{11})$ is a self-adjoint contraction, define
\[
  \Theta_W=\arccos\!\left(\frac{A_\perp}{2\sqrt{11}}\right).
\]

\begin{theorem}[Self-adjoint graph-RH phase determinant]
The operator $\Theta_W$ is self-adjoint and
\[
 \det\!\left(
 I-2\sqrt{11}\,u\cos\Theta_W+11u^2I
 \right)
 =(1-2u+11u^2)^{24}(1+4u+11u^2)^{15}.
\]
Consequently the nontrivial Ihara poles lie on $|u|=11^{-1/2}$, and under
$u=11^{-s}$ their logarithmic lifts lie on $\Re(s)=1/2$.
\end{theorem}

\section{Exact pole arithmetic}

The four distinct nontrivial pole locations are
\[
  \frac{1\pm\sqrt{-10}}{11}
  \quad\text{of order }24,
  \qquad
  \frac{-2\pm\sqrt{-7}}{11}
  \quad\text{of order }15.
\]
Thus the restricted multiplicity count $24+15=39$ corresponds to $39$
adjacency slots, four distinct pole locations, and total nontrivial pole order
$2\cdot24+2\cdot15=78$.

The coordinate denominator ideals are generated by the conjugate norm-$11$
elements
\[
  1\mp\sqrt{-10},
  \qquad
  -2\mp\sqrt{-7},
\]
respectively.  All principal-part coefficients and residues are computed
exactly by \texttt{analysis/w33\_ihara\_principal\_parts.py}.

\section{Finite countermodel classification}

Put $z=s-\tfrac12$.  A real polynomial with reflection symmetry
$F(1-s)=F(s)$ has the form
\[
  F(s)=Q(z^2),\qquad Q\in\mathbb R[y].
\]
Its zeros lie on $\Re(s)=1/2$ if and only if all roots of $Q$ are real and
nonpositive.  Equivalently, up to a real constant,
\[
  F(s)=\det(z^2I+H^2)
\]
for a finite self-adjoint operator $H$.

For one reflected quartet,
\[
 F=((z-\delta)^2+\gamma^2)((z+\delta)^2+\gamma^2)
\]
and
\[
 Q(y)=y^2+2(\gamma^2-\delta^2)y+(\gamma^2+\delta^2)^2,
 \qquad
 \operatorname{disc}Q=-16\gamma^2\delta^2.
\]
For $\gamma\ne0$, the self-adjoint-square factorization holds exactly when
$\delta=0$.

\section{Classical transfer obstruction}

The present exact chain is
\[
 \text{completed-}\xi\text{ current}
 \longrightarrow
 \text{reflection cocycle}
 \longrightarrow
 \text{positive orbit energy}
 \longrightarrow
 \text{exact zero-defect }W(3,3)\text{ model}.
\]
A classical proof still requires one of the following:

\begin{enumerate}
  \item a trace formula identifying a convergent total cocycle energy and
        proving that it is zero;
  \item a Weil/Li/de Branges-type positive Hermitian form shown to coincide
        with the defect functional;
  \item an infinite self-adjoint $W(3,3)$ tower whose regularized determinant
        is exactly the completed Riemann $\xi$ function.
\end{enumerate}

The finite $W(3,3)$ phase operator is exact for Ihara zeta.  Direct comparison
of the quadratic, quartic, and sextic logarithmic Hadamard coefficients shows
that neither its principal spectrum nor its rigid logarithmic tower currently
matches classical $\xi$ under a single scale.
